\documentclass[11pt]{article}

\usepackage[T1]{fontenc}
\usepackage{lmodern}
\usepackage{amsmath,amssymb,amsthm}
\usepackage[margin=1in]{geometry}
\usepackage{microtype}
\usepackage{needspace}
\usepackage[hidelinks]{hyperref}
\hypersetup{
  pdftitle={Dense critical graphs from saturated graphs},
  pdfauthor={Qiyuan Gu}
}

\newtheorem{theorem}{Theorem}
\newtheorem{lemma}[theorem]{Lemma}
\newtheorem{proposition}[theorem]{Proposition}
\theoremstyle{remark}
\newtheorem*{remark}{Remark}
\numberwithin{equation}{section}
\newcommand{\col}{\operatorname{col}}

\title{Dense critical graphs from saturated graphs}
\author{Qiyuan Gu\\
University of Chicago\\
\href{mailto:phoenix1203@uchicago.edu}{\texttt{phoenix1203@uchicago.edu}}}
\date{}

\begin{document}
\maketitle

\begin{abstract}
Let \(f_k(n)\) denote the maximum number of edges in a \(k\)-critical
graph on \(n\) vertices.
For every fixed integer \(k\ge8\), we construct a sequence of
\(k\)-critical graphs \(G_s\) with \(|V(G_s)|\to\infty\) and
\[
 \frac{e(G_s)}{|V(G_s)|^2}\longrightarrow
 \begin{cases}
 \dfrac{k-4}{2(k-2)},&k\text{ even},\\[4pt]
 \dfrac{k-5}{2(k-3)},&k\text{ odd}.
 \end{cases}
\]
For \(k=12\), the density is \(2/5>3/8\), disproving
\(f_{12}(n)\sim3n^2/8\) and hence the general asymptotic formula
proposed in Erd\H{o}s Problem 917.
The construction combines a coloring module of Pegden with
saturated graphs of sublinear maximum degree.
\end{abstract}

\section{Introduction}

All graphs are finite and simple. A graph is \(k\)-critical if it has
chromatic number \(k\) and every proper subgraph is \((k-1)\)-colorable;
for graphs without isolated vertices this is equivalent to requiring
only that deleting any edge lowers the chromatic number.
Write \(e(G)=|E(G)|\). For each integer \(n\ge0\), let \(f_k(n)\) be
the largest edge count of a \(k\)-chromatic graph on \(n\) vertices
satisfying this edge-deletion condition, with \(f_k(n)=0\) if no such
graph exists.
Erd\H{o}s \cite[p.~28]{Erdos1969} conjectured that, for fixed \(k\ge6\),
\begin{equation}\label{eq:conjecture}
 f_k(n)\sim\frac12\left(1-\frac1{\lfloor k/3\rfloor}\right)n^2.
\end{equation}
This is the general asymptotic question in Erd\H{o}s Problem 917
\cite{Erdos917}.
Dirac \cite[p.~90]{Dirac1952} gave the \(k=6\) construction obtained by
joining two odd cycles of equal order, and Erd\H{o}s
\cite[p.~28, eq.~(3)]{Erdos1969} observed that it generalizes when
\(3\mid k\), giving the coefficient in \eqref{eq:conjecture}.
When \(3\nmid k\), Toft's constructions \cite{Toft1970} give larger
coefficients. For background on dense critical graphs, see Jensen
and Toft \cite[Section~5.1, pp.~97--98]{JensenToft} and Jensen
\cite{Jensen2002}. Luo, Ma, and Yang \cite{LMY} noted that no
constructions improving Toft's asymptotic constants had been found
since 1970. We give the following construction.

\begin{theorem}\label{thm:main}
For every integer \(k\ge8\), there is a sequence of \(k\)-critical
graphs \(G_s\) such that \(|V(G_s)|\to\infty\) and
\[
 \frac{e(G_s)}{|V(G_s)|^2}\longrightarrow
 \begin{cases}
 \dfrac{k-4}{2(k-2)},&k\text{ even},\\[4pt]
 \dfrac{k-5}{2(k-3)},&k\text{ odd}.
 \end{cases}
\]
\end{theorem}

For even \(k=2t+2\), the graph \(G\) consists of \(t\) copies of
Pegden's module \(U(S,T)\) \cite{Pegden}, with their active sets
joined by an almost complete \(t\)-partite graph. The omitted edges
form a graph \(Z\) whose adjacency is prescribed by a
\(K_t\)-saturated graph \(H\) of small maximum degree from \cite{AEHK}.
In a \((2t+1)\)-coloring, each active set needs three colors
(Lemma~\ref{lem:module}); the degree bound forces two of them to be
private to that part. The remaining color must then appear in all
\(t\) parts, giving a \(K_t\) in \(Z\), impossible since \(H\) is
\(K_t\)-free. Deleting an edge between two parts adds an edge to
\(Z\), and saturation supplies a \(K_t\) through it, allowing one
color to be shared in a \((2t+1)\)-coloring of \(G-e\).
Complete joins of critical graphs, as in Dirac's construction,
cannot share colors between factors; omitting \(Z\) costs \(o(n^2)\)
edges and permits this shared color.

Deleting the \(O_k(n)\) module edges leaves a \(t\)-partite graph,
so the asymptotic density of this construction is at most
\((t-1)/(2t)\), equal to \(2/5\) when \(k=12\).
For comparison, the upper bound of Luo, Ma, and Yang
\cite[Theorem~1.1]{LMY} has coefficient
\(9/20-1/4356\approx0.44977\) for \(f_{12}\).
Section~\ref{sec:parameters} gives an explicit finite-field
realization of the \(k=12\) case.

\section{A coloring module}
\label{sec:module}

Let \(r\ge3\), let \(S\) be \(r\)-critical, and let \(T\) be
\((r-1)\)-critical.
The module \(U(S,T)\) consists of disjoint copies of \(S\) and \(T\),
together with an independent set
\[
 A=V(S)\times V(T).
\]
Each \((x,y)\in A\) is adjacent to \(x\in V(S)\) and
\(y\in V(T)\). There are no edges between \(S\) and \(T\), and
no other edges. We call \(A\) the active set and call
\(\{x\}\times V(T)\) the row indexed by \(x\).

The following is the \(U(r,r-1)\) case of Pegden's module lemma
\cite[Lemma~2.5]{Pegden}, stated there for triangle-free factors;
the lemma holds for general critical factors \(S\) and \(T\) of
the indicated chromatic numbers.

\begin{lemma}\label{lem:module}
Fix a palette of \(r\) colors.
\begin{enumerate}
\item In every proper coloring of \(U(S,T)\) from this palette,
the active set uses at least three colors.
\item Given \(a_0=(x_0,y_0)\in A\) and distinct colors
\(\alpha,\beta,\gamma\), there is a proper coloring in which
the active set uses only \(\alpha,\beta,\gamma\), and
\(\gamma\) occurs on the active set only at \(a_0\).
\item Given an edge \(e\) of \(U(S,T)\) and distinct colors
\(\alpha,\beta\), there is a proper coloring of \(U(S,T)-e\)
in which the active set uses only \(\alpha,\beta\).
\end{enumerate}
\end{lemma}

\begin{proof}
For (1), suppose the active colors are contained in
\(\{\alpha,\beta\}\). Since \(S\) uses all \(r\) colors, it
contains a vertex colored \(\alpha\) and a vertex colored
\(\beta\). Every active vertex in the first row must have color
\(\beta\), so no vertex of \(T\) can have color \(\beta\).
The second row similarly excludes \(\alpha\) from \(T\).
This leaves only \(r-2\) colors for \(T\), a contradiction.

For (2), color \(S\) so that \(\alpha\) occurs only at \(x_0\).
Color \(T\) with the \(r-1\) colors other than \(\alpha\), so that
\(\beta\) occurs only at \(y_0\). These colorings exist by criticality:
delete the specified vertex, color the remaining graph with one
fewer color, and give the deleted vertex its own color.
Extend these colorings by
\begin{equation}\label{eq:singleton-coloring}
 \col(x,y)=
 \begin{cases}
 \alpha,&x\ne x_0,\\
 \beta,&x=x_0,\ y\ne y_0,\\
 \gamma,&(x,y)=(x_0,y_0).
 \end{cases}
\end{equation}
Each active vertex differs in color from both its neighbors.

For (3), there are four possibilities for the deleted edge.
If \(e\in E(S)\), color both \(S-e\) and \(T\) with the \(r-1\)
colors other than \(\alpha\), and color all of \(A\) with
\(\alpha\). If \(e\in E(T)\), color \(T-e\) with the \(r-2\)
colors other than \(\alpha,\beta\). Choose \(x_0\in V(S)\)
and color \(S\) with \(\alpha\) only at \(x_0\). Color the
row indexed by \(x_0\) with \(\beta\) and all other rows with
\(\alpha\).

Suppose instead that \(e\) joins an active vertex
\((x_0,y_0)\) to one of its two structural neighbors.
Use the colorings of \(S\) and \(T\) from the proof of (2).
If the deleted edge joins \((x_0,y_0)\) to \(x_0\), use
\eqref{eq:singleton-coloring} with the color of \((x_0,y_0)\)
changed to \(\alpha\). If it joins \((x_0,y_0)\) to \(y_0\),
change that color to \(\beta\). In either case the only edge
that would become monochromatic is the deleted edge.
\end{proof}

\section{From saturated graphs to critical graphs}
\label{sec:conversion}

A graph is \(K_t\)-saturated if it contains no \(K_t\), but adding
any missing edge creates a \(K_t\). Thus the common neighborhood
of every pair of distinct nonadjacent vertices contains a \(K_{t-2}\).

\begin{proposition}\label{prop:conversion}
Let \(t\ge3\), and let \(H\) be a \(K_t\)-saturated graph on
\(v\ge t\) vertices with maximum degree \(d<v-1\).
Let \(h\ge2t+1\) be odd and suppose
\begin{equation}\label{eq:degree-condition}
 v>2t(t-1)hd.
\end{equation}
Put \(a=2hv\). There is a \((2t+2)\)-critical graph \(G\) with
\begin{equation}\label{eq:order}
 |V(G)|=t(a+h+2v)
\end{equation}
and
\begin{equation}\label{eq:edge-bound}
 e(G)\ge\binom t2 a^2\left(1-\frac d v\right).
\end{equation}
\end{proposition}

\begin{proof}
\emph{Construction.}
Set \(\ell=2v\) and take joins of cliques and odd cycles, as in
\cite[p.~90]{Dirac1952}:
\[
 S=K_{2t-2}\vee C_{h-2t+2},\qquad
 T=K_{2t-3}\vee C_{\ell-2t+3}.
\]
Both cycles are odd and have length at least three. Since the join of
a \(p\)-critical and a \(q\)-critical graph is \((p+q)\)-critical,
\(S\) is \((2t+1)\)-critical and \(T\) is \(2t\)-critical.
(Deleting a vertex or an edge within one factor saves a color there;
if a joining edge \(xy\) is deleted, give \(x\) and \(y\) singleton
colors in their respective factors and identify those colors.)

Take \(t\) disjoint copies \(U_1,\ldots,U_t\) of \(U(S,T)\),
with active sets \(A_1,\ldots,A_t\), each of size \(a=h\ell\).
In each copy, identify \(V(T)\) with \(V(H)\times\{0,1\}\)
by any bijection. An active vertex then has the form
\((x,(z,\varepsilon))\); give it label \(z\).
Let \(\pi\) denote this labeling map. Each label occurs
exactly \(b=2h\) times in each active set.

Define a \(t\)-partite graph \(Z\) on \(A_1\cup\cdots\cup A_t\)
by putting, for \(u\in A_i\) and \(w\in A_j\) with \(i\ne j\),
\[
 uw\in E(Z)\quad\Longleftrightarrow\quad
 \pi(u)\pi(w)\in E(H).
\]
The graph \(G\) retains all the edges of the modules and,
between different active sets, takes the complement of \(Z\).
There are no other edges between modules. In particular, active
vertices with equal labels in different parts are adjacent in \(G\).

We record three properties of \(Z\). First, \(Z\) is
\(K_t\)-free: the labeling map sends every clique injectively
to a clique of \(H\). Moreover, adding any missing edge between
different parts of \(Z\) creates a \(K_t\) with one vertex in
each part. To see this, let \(u,w\) be the endpoints of such a
missing edge. If their labels \(z,z'\) differ, saturation of
\(H\) gives a \(K_{t-2}\) in \(N_H(z)\cap N_H(z')\).
Place its labels in the other \(t-2\) parts. Together with
\(u,w\), they form the required clique in \(Z+uw\).
If the labels coincide, say both equal \(z\), choose a vertex
\(z'\ne z\) not adjacent to \(z\); this is possible because
\(d<v-1\). The same common-neighborhood clique is contained
in \(N_H(z)\) and gives the required clique.

Second, \(Z\) contains a \(K_{t-1}\) across any prescribed
\(t-1\) parts. Indeed, \(H\) has a missing edge; one endpoint
together with a \(K_{t-2}\) in the common neighborhood gives
a \(K_{t-1}\) in \(H\). Its labels can be placed in any
\(t-1\) parts of \(Z\).

Third, \(D:=\Delta(Z)=(t-1)bd=2h(t-1)d\), so
\eqref{eq:degree-condition} gives \(\ell>2tD\).

\smallskip
\noindent\emph{The chromatic lower bound.}
Suppose that \(G\) has a \((2t+1)\)-coloring.
Fix a part \(A_i\) and any color \(\alpha\).
Since its copy of \(S\) uses all \(2t+1\) colors, some vertex
\(x\) of \(S\) has color \(\alpha\). The \(\ell\) active
vertices in the row indexed by \(x\) avoid \(\alpha\), so the
pigeonhole principle and \(\ell>2tD\) give a color
\(\beta\ne\alpha\) occurring more than \(D\) times in that row.
Applying this observation twice, we may choose two distinct
colors \(\alpha_i,\beta_i\), each used more than \(D\) times in \(A_i\).

A color used more than \(D\) times in \(A_i\) cannot occur in
any other active set. A vertex of that color in another part
would have to be adjacent in \(Z\) to all those more than
\(D\) vertices, contrary to \(\Delta(Z)=D\).
Thus the \(2t\) chosen colors
\(\alpha_1,\beta_1,\ldots,\alpha_t,\beta_t\) are distinct,
leaving a single color \(\gamma\).
The active set \(A_i\) can use only
\(\alpha_i,\beta_i,\gamma\), and Lemma~\ref{lem:module}(1)
forces \(\gamma\) to occur in every part.
Choosing a vertex of this color in each part gives a
\(K_t\) in \(Z\), a contradiction. Thus \(\chi(G)\ge2t+2\).

\smallskip
\noindent\emph{Colorings after edge deletion.}
Use the \((2t+1)\)-color palette
\[
 \alpha_1,\beta_1,\ldots,\alpha_t,\beta_t,\gamma.
\]
The pair \(\alpha_i,\beta_i\) will be used only in \(A_i\)
among the active sets. Structural vertices may use the full
palette, since they have no neighbors in other modules.

Let \(e=uw\) join two different active sets in \(G\).
By the first property of \(Z\), there is a \(K_t\) in
\(Z+uw\) containing \(uw\), with vertices \(a_i\in A_i\).
Apply Lemma~\ref{lem:module}(2) in each module, using active
colors \(\alpha_i,\beta_i,\gamma\) and assigning \(\gamma\)
only to \(a_i\) in \(A_i\). The only color shared by different
active sets is \(\gamma\). Its \(t\) vertices are pairwise
nonadjacent in \(G-e\), because they form a clique in \(Z+uw\).
We have obtained a \((2t+1)\)-coloring of \(G-e\).

Now let \(e\) lie within \(U_i\). Apply
Lemma~\ref{lem:module}(3) to color \(U_i-e\), using only
\(\alpha_i,\beta_i\) on \(A_i\). Choose a \(K_{t-1}\) in
\(Z\) across the other \(t-1\) parts. In each of those modules,
apply Lemma~\ref{lem:module}(2) with the selected vertex as
the only active vertex of color \(\gamma\). Once again this
is a proper \((2t+1)\)-coloring of \(G-e\).

These cases cover every edge of \(G\). Recoloring one endpoint
of a deleted edge with a new color gives \(\chi(G)\le2t+2\).
Since \(G\) has no isolated vertices, every proper subgraph of
\(G\) lies in some \(G-e\). Hence \(G\) is \((2t+2)\)-critical.

\smallskip
\noindent\emph{Counting vertices and edges.}
Each module has \(a+h+2v\) vertices, giving \eqref{eq:order}.
Between a fixed pair of active sets, \(Z\) has \(2b^2m\) edges:
each edge of \(H\) gives two ordered pairs of labels, with
\(b^2\) pairs of active vertices for each, where \(m=e(H)\).
Since \(2b^2m\le b^2vd\) and \(a=bv\), at least
\(a^2(1-d/v)\) edges remain between each pair of active sets,
proving \eqref{eq:edge-bound}.
\end{proof}

\begin{remark}[Exact edge count]
The construction in Proposition~\ref{prop:conversion} satisfies
\begin{equation}\label{eq:exact-edges}
 e(G)=\binom t2(a^2-8h^2m)
 +t\bigl(2a+(2t-1)h+4(t-1)v-(4t^2-4t-1)\bigr),
\end{equation}
where \(m=e(H)\).
Indeed,
\[
 e(S)=(2t-1)h-(t-1)(2t+1),\qquad
 e(T)=4(t-1)v-t(2t-3),
\]
and each module has \(2a\) edges incident with its active set.
\end{remark}

\begin{proof}[Proof of Theorem~\ref{thm:main}]
First let \(k=2t+2\) be even, where \(t\ge3\) is fixed.
By \cite[Theorem~7]{AEHK}, there are \(K_t\)-saturated graphs
\(H\) of order \(v\to\infty\) and maximum degree
\(d=O_t(\sqrt v)\). Choose
\[
 h=2\lfloor v^{1/4}\rfloor+1.
\]
Then \(h\to\infty\), \(h=o(\sqrt v)\), and
\(2t(t-1)hd/v\to0\). Consequently \(h\ge2t+1\), \(v\ge t\),
\(d<v-1\), and \eqref{eq:degree-condition} hold for all
sufficiently large \(v\).

Apply Proposition~\ref{prop:conversion}, and write
\(a=2hv\) and \(N=|V(G)|\). We have
\[
 \frac{N}{ta}=1+\frac1h+\frac1{2v}\longrightarrow1,
 \qquad \frac d v\longrightarrow0.
\]
The edges within modules number \(O_t(a)\), by
\eqref{eq:exact-edges}, and hence
\[
 \binom t2a^2\left(1-\frac d v\right)
 \le e(G)\le\binom t2a^2+O_t(a).
\]
It follows that
\[
 \frac{e(G)}{N^2}\longrightarrow
 \frac{\binom t2}{t^2}=\frac{t-1}{2t}
 =\frac{k-4}{2(k-2)}.
\]
For odd \(k=2t+3\ge9\), join each graph in the even construction
to \(K_1\). The resulting graph is \(k\)-critical, has \(N+1\)
vertices and \(e(G)+N\) edges, and has the same limiting density,
now equal to \((k-5)/(2(k-3))\).
\end{proof}

\section{An explicit family for \texorpdfstring{\(k=12\)}{k = 12}}
\label{sec:parameters}

We use the \(k=5\) case of \cite[Section~4, Example~2]{AEHK} and
include the saturation verification in Lemma~\ref{lem:saturation}.

\begin{lemma}[Alon--Erd\H{o}s--Holzman--Krivelevich]
\label{lem:saturation}
For every prime power \(Q\ge4\), there is a \(K_5\)-saturated
graph \(H_Q\) with
\begin{equation}\label{eq:H-parameters}
 |V(H_Q)|=13Q^2+12Q,
 \qquad \Delta(H_Q)=22Q-3.
\end{equation}
\end{lemma}

\begin{proof}
We give the verification of Example 2 in \cite{AEHK} for \(k=5\).
Let \(F\) be a field of order \(Q\). The levels are
\(F\cup\{\infty\}\), given any linear order, and the places in
each level are indexed by \(\mathbb Z/Q\mathbb Z\) through a fixed
bijection with \(F\). For each \((a,b)\in F^2\), form a line
\(L_{a,b}\) whose point in finite level \(u\) has field coordinate
\(au+b\), and whose point in level \(\infty\) has coordinate \(a\).
Every line has one point in each level, every point lies on \(Q\)
lines, and any two distinct lines meet: if their slopes agree they
meet in level \(\infty\); otherwise solve
\((a-a')u=b'-b\).

A core vertex is a quadruple \((i,j,t,c)\), where \(i\) is its
level, \(j\in\mathbb Z/Q\mathbb Z\) its place,
\(t\in\mathbb Z/3\mathbb Z\) its type, and
\(c\in\{0,1,2,3\}\) its copy. In one level, two core vertices
are adjacent precisely when their copies differ and their
\((j,t)\) pairs differ. Between levels \(i<i'\), adjacency means
\[
 t'=t+1,\qquad j'\in\{j+1,j+2,j+3\}.
\]
Finally, add one vertex for each line, adjacent to all twelve core
vertices at each of its points. The line vertices are independent.
All additions in the place coordinate are modulo \(Q\), irrespective
of the characteristic of \(F\).

\emph{Excluding \(K_5\).}
A triangle of core vertices cannot occupy three levels: its edges
would require both \(t''=t+1\) and \(t''=t+2\) modulo three.
Thus every core clique occupies at most two levels. Within one
level it has at most four vertices, since copies must differ.
If it meets two levels, all vertices in each level have the same
type, so their places and copies are distinct. The three successors
of a place form a proper cyclic interval, since \(Q\ge4\), and
different places give different successor triples. The same is true
of predecessor triples. Hence two distinct places have at most two
common successors or predecessors. A clique meeting two levels
therefore cannot have three vertices in one level and two in the
other, nor four in either level.

The neighborhood of a line vertex is \(K_4\)-free. Indeed, a clique
in that neighborhood occupying one level has at most three vertices,
as its places agree and its types must differ. If it occupies two
levels, the preceding same-type observation permits at most one
vertex in each. Since line vertices are independent, this also
excludes every \(K_5\) containing a line vertex.

\emph{Saturating every missing edge.}
It suffices to exhibit a triangle in the common neighborhood of
any two distinct nonadjacent vertices \(x,y\).
If both are line vertices, take their intersection point and three
core vertices there with distinct types and copies.
If \(x\) is a core vertex and \(y\) a nonincident line, take the
point of \(y\) in the level of \(x\). Its place differs from that
of \(x\); use all three types there, in the three copies other
than the copy of \(x\).

Suppose next that both vertices are in the same level. If their
copies agree, choose a place different from both of their places,
and use its three types in the three other copies. Such a place
exists because \(Q\ge4\). If their copies differ, nonadjacency
means that their places and types agree. Choose a line through
that point and two core vertices there using the other two types
and the other two copies. These three vertices give the triangle.

Finally, suppose \(x\) lies below \(y=(i',j',t',c')\).
In the level of \(x\), take the three places
\(j'-1,j'-2,j'-3\), all of type \(t'-1\), assigning them the
three copies different from the copy of \(x\). They form a triangle
and are adjacent to \(y\). None has the same place and type as
\(x\), since that would make \(x\) adjacent to \(y\).
Thus they are also adjacent to \(x\). This exhausts the missing edges.

\emph{Order and maximum degree.}
There are \(12Q(Q+1)\) core vertices and \(Q^2\) line vertices.
A core vertex has \(3(3Q-1)\) neighbors in its level,
\(12Q\) in the other levels, and \(Q\) line neighbors.
Its degree is therefore \(22Q-3\). A line vertex has degree
\(12(Q+1)<22Q-3\) for \(Q\ge4\), proving
\eqref{eq:H-parameters}. In particular,
\[
 \frac{\Delta(H_Q)}{|V(H_Q)|}
 \le \frac{22}{13Q}\longrightarrow0
 \qquad(Q\longrightarrow\infty).
\]
\end{proof}

\begin{proof}[The case \(k=12\)]
For each integer \(s\ge4\), set
\[
 h=3^s,\qquad Q=h^2,
\]
and take \(H=H_Q\) from Lemma~\ref{lem:saturation}. Write
\(v=|V(H)|\) and \(d=\Delta(H)\). Then
\[
 v=13h^4+12h^2,\qquad d=22h^2-3<v-1.
\]
The integer \(h\) is odd and at least \(81\). Moreover,
\[
 40hd<880h^3<13h^4\le v,
\]
since \(13h\ge13\cdot81>880\).
Proposition~\ref{prop:conversion} with \(t=5\) gives a twelve-critical
graph \(G_s\) of order \(N_s=5(2hv+h+2v)\).
Since \(N_s/(10hv)\to1\) and \(d/v\to0\),
\eqref{eq:edge-bound} and \eqref{eq:exact-edges} give
\(e(G_s)/N_s^2\to2/5\).
Finally, \(N_s\to\infty\) and \(f_{12}(N_s)\ge e(G_s)\), so
\[
 \limsup_{n\to\infty}\frac{f_{12}(n)}{n^2}\ge\frac25>\frac38.
\]
Thus \(f_{12}(n)\not\sim3n^2/8\).
\end{proof}

\section*{Code availability}

A Lean 4 formalization of the \(k=12\) case of Theorem~\ref{thm:main},
Lemma~\ref{lem:saturation}, and the corresponding instances of
Lemma~\ref{lem:module} and Proposition~\ref{prop:conversion} is available at
\url{https://github.com/FireflySentinel/erdos-917}.
The general result for other chromatic numbers is not formalized.

\section*{Declaration of generative AI and AI-assisted technologies in the writing process}

GPT-6 Astra proposed the construction of critical graphs and the
density estimates, drafted the manuscript, and generated the Lean formalization.
GPT-5.6 Sol and Claude Opus 5 were used to organize earlier research material
and to review the Lean code.
The author checked the arguments, completed the final manuscript, and takes
full responsibility for the content.

\begin{thebibliography}{9}

\bibitem{AEHK}
N.~Alon, P.~Erd\H{o}s, R.~Holzman, and M.~Krivelevich,
On \(k\)-saturated graphs with restrictions on the degrees,
\emph{J. Graph Theory} \textbf{23} (1996), 1--20.
\url{https://web.math.princeton.edu/~nalon/PDFS/aehk1.pdf}.

\bibitem{Dirac1952}
G.~A. Dirac,
A property of 4-chromatic graphs and some remarks on critical graphs,
\emph{J. London Math. Soc.} \textbf{27} (1952), 85--92.
\href{https://doi.org/10.1112/jlms/s1-27.1.85}{doi:10.1112/jlms/s1-27.1.85}.

\bibitem{Erdos1969}
P.~Erd\H{o}s,
Problems and results in chromatic graph theory,
in \emph{Proof Techniques in Graph Theory}
(Proc. Second Ann Arbor Graph Theory Conf., Ann Arbor, 1968),
Academic Press, New York, 1969, 27--35.
\url{https://www.renyi.hu/~p_erdos/1969-13.pdf}.

\bibitem{Erdos917}
T.~F. Bloom,
\emph{Erd\H{o}s Problem \#917},
\url{https://www.erdosproblems.com/917}.
Accessed 5 September 2026.

\bibitem{Jensen2002}
T.~R. Jensen,
Dense critical and vertex-critical graphs,
\emph{Discrete Math.} \textbf{258} (2002), 63--84.
\href{https://doi.org/10.1016/S0012-365X(02)00262-5}{doi:10.1016/S0012-365X(02)00262-5}.

\bibitem{JensenToft}
T.~R. Jensen and B.~Toft,
\emph{Graph Coloring Problems},
Wiley-Interscience, New York, 1995.
\href{https://doi.org/10.1002/9781118032497.ch5}{Chapter~5: Critical graphs}.

\bibitem{LMY}
C.~Luo, J.~Ma, and T.~Yang,
On the maximum number of edges in \(k\)-critical graphs,
\emph{Combin. Probab. Comput.} \textbf{32} (2023), 900--911.
\href{https://doi.org/10.1017/S0963548323000238}{doi:10.1017/S0963548323000238}.

\bibitem{Toft1970}
B.~Toft,
On the maximal number of edges of critical \(k\)-chromatic graphs,
\emph{Studia Sci. Math. Hungar.} \textbf{5} (1970), 461--470.

\bibitem{Pegden}
W.~Pegden,
Critical graphs without triangles: an optimum density construction,
\emph{Combinatorica} \textbf{33} (2013), 495--512.
\href{https://arxiv.org/abs/1101.4417v2}{arXiv:1101.4417v2}.

\end{thebibliography}

\end{document}
